\section{Conclusion}
In this work, we presented a conceptually simple and effective two-stage training schema to learn bilingual and multilingual multimodal representation models, via teacher learning and contrastive learning. We show the effectiveness of our method by conducting extensive experiments on a wide range of tasks in English and Chinese. In Chinese, our method sets new state-of-the-art results on multiple zero-shot image classification and retrieval tasks while being highly effective: we only use tens of millions of text data and two million text-image pairs during training, whereas most prior work requires training on hundreds of millions of text-image pairs. Future work includes exploring altering the image encoder to combine vision signals learned from different data distributions and possibly eliminating the need for machine-translated data to build a multilingual multimodal pretraining model.

%Learning a good representation in a joint space for vision and language has been a long pursuit in the research of Artificial Intelligence (AI). Recently, the milestone work of CLIP \cite{clip} from OpenAI demonstrated impressive zero-shot performances across a number of tasks such as Image classification on ImageNet, Image-to-Text and Text-to-Image retrieval on Flicker-30K, MSCOCO, . There has been pursuit of contrastive language-image model in other languages such as Italian, Korean, Chinese or in a cross-lingual/multi-lingual setting \cite{}.

%Training a language-image representation usually require a large number of text-image pairs (for instance, CLIP used 400M text-image pairs, and Taiyi, a Chinese CLIP model, used 123M text-image pairs), and a huge number of computing resources where most research groups don’t have access to.  \cite{} Proposed a novel technique to Teacher Learning (a.k.a. Knowledge Distillation) on the textual encoder of the CLIP model, which only require machine-translated data from English to target language, without text-to-image pairs. We follow this line of work and further improved their method in the following ways.

%We proposed a simple, scalable, and effective two stage approach for learning a good Chinese CLIP model while keeping the same capability.
%Notably, we achieved on-par performance on image retrieval tasks for both English and Chinese. However, the results on Imagnet differs a lot: while CLIP achieves xxx on ImageNet~\cite{deng2009imagenet} classification, our model drops at an accuracy of 59.xx\% on Imagnet-cn. We suspect this is due to xxx. We left it as future direction to close this gap on Imagnet between English and Chinese.

